\thispagestyle{plain}
\begin{center}
{\fontsize{10}{12}\selectfont Manuscrit de compilation pour le transfert académique\par}
\vspace{1.5mm}
{\fontsize{17}{20.5}\selectfont\bfseries Relations structurelles\par entre les constantes physiques\par}
\vspace{1.5mm}
{\fontsize{11}{13.5}\selectfont Action, réciprocité géométrique\par et mémoire des transformations\par}
\vspace{1.5mm}
{\fontsize{11.5}{14}\selectfont Oumar Haidara Fall \textperiodcentered\ Rubén Ramos Balsa\par}
\vspace{1.5mm}
{\fontsize{8}{10}\selectfont Coautorat sans préséance de contribution\par}
\vspace{1.5mm}
{\fontsize{8}{10}\selectfont Intégration et compilation documentaires générées par ChatGPT - Astra.\par}
\end{center}
\vspace{1.5mm}
\begingroup\fontsize{9.5}{10.7}\selectfont\setlength{\parskip}{1.5pt}\renewenvironment{abstract}{\begin{center}\bfseries\abstractname\end{center}\noindent}{\par}\begin{abstract}
On établit des relations de compatibilité, de réciprocité et de reconstruction
entre les structures d’action, de masse et de réponse constitutive obtenues
par l’holographie modulaire triadique et la mécanique dimensionnelle. Le développement
part de la composition de la structure arithmétique des chemins
\emph{All Possible Paths} (APP), de la signature ternaire orientée TRIT et du
\emph{Topological Phase Kernel} (TPK), qui met à jour et transporte l’état
enrichi. Cette composition conserve orientation, retenue, mémoire et
retour. Les constantes interviennent comme évaluations des
structures préalablement construites ; leur composition permet de déterminer quelles
propriétés elles partagent et quelle information une observation conjointe récupère.
Dans la collection constitutive normalisée, on démontre que la vitesse de
propagation et la fonctionnelle de Barbero--Immirzi déterminent les valeurs propres
étiquetées des deux canaux. Les logarithmes de vitesse et d’impédance
constituent le gradient d’un potentiel spectral strictement concave.
Sa hessienne caractérise la sensibilité de l’inversion et fournit des bornes
explicites de stabilité. Le potentiel partage un moment spectral de
profondeur deux avec la construction orientée dont la valeur extrême est Catalan ;
la réponse constitutive restitue cette collection de manière unique avec sa
condition de frontière à l’origine.

La composition de quarts de tour associés à des formes elliptiques conjuguées
produit un transport relatif unimodulaire dont les dilatations réciproques
permettent de reconstruire le coefficient adimensionnel d’action, la
coordonnée de clôture et l’orientation étant conservées. Le relèvement dodécaphasique et
l’itération du transport caractérisent la mémoire de l’ordre des
opérations. L’analyse chronologique distingue la croissance de
l’information du taux entropique par position, et la réentrée de mémoire
détermine sa contribution à la réponse opératoire. Le retour inscrit et
sa composition polaire déterminent une longueur $L_\alpha$. Conserver
conjointement action et aire fixe $G=c^3L_\alpha^2/\hbar$ dans la réalisation
radiale, également sous des variations non commutatives et une réduction compatible
de mémoire.

La mémoire électronique est restituée par des invariants linéaires et
spectraux qui conservent l’ordre du registre ; le transport entre
multisections admet un critère explicite de conservation de masse.
La composition thermique réunit capacités, multiplicités et mémoire
spectrale dans une trace exacte, et construit sa réponse thermométrique
pour une lecture marquée à référence fixe. Cette réalisation relie
énergie, information, action et aire, en conservant les conditions de
normalisation et d’observation de chaque grandeur.
La composition transportée des réponses et la reconstruction
Higgs--impédance relient les canaux constitutifs au secteur
électrofaible. La substitution électronique conserve son échelle dans Fermi,
la simplifie dans les quotients angulaires et transporte la longueur inverse
par une congruence opératoire, également sans commutation.

Enfin, sous les
conditions circulaire et thermique spécifiées, une même coordonnée de masse
détermine la fréquence propre, deux longueurs réciproques de produit
invariant et la pondération thermique du cycle. Les résultats réunissent ainsi
production, compatibilité et restitution dans une description mathématique
commune, avec des démonstrations explicites et la distinction entre changement de
représentation, transport relatif et réalisation physique.


Le lecteur dodécaphasique uniforme détermine également une suite de premières racines de retour critique. On démontre que le rapport de leurs écarts converge vers la constante de Feigenbaum, en conservant la sélection initiale, l’orientation et la mémoire de la construction.

% Adición al resumen de VII; no sustituye su contenido anterior.
La même composition se prolonge en une action totale de géométrie, de jauge et de matière : elle incorpore conjointement la gravitation et les interactions forte, faible et électromagnétique, en conservant les représentations chirales et les couplages précédemment générés. On obtient son courant de spin, la source métrique totale et les termes torsionnels croisés entre espèces. La transformation de Legendre et les identités de symétrie produisent des contraintes conjointes de première classe et leur propagation. Sur leur surface régulière, on calcule la courbure de la connexion canonique totale et l’on démontre son annulation dans les directions de déformation caractéristiques. Le résultat conserve les conditions de bord et l’éventuelle holonomie globale de l’état enrichi.



\par\smallskip\noindent\textbf{Mots-clés :} action ; réciprocité géométrique ; mémoire opératoire ; constantes physiques ; holographie modulaire triadique ; mécanique dimensionnelle.
\end{abstract}


\endgroup
